\chapter{Neuronale Netze}
\label{chap:neuronale-netze}

\begin{myremark}[Algorithmusart in diesem Kapitel]
Funktionsart: \textbf{Klassifikation und Regression}. \\
Umsetzungsart: \textbf{datengetrieben}.
\end{myremark}

\begin{myremark}
Dieses Kapitel ist anspruchsvoller als die übrigen.
Sie setzt Kapitel~\cref{chap:einfuehrung} und Kapitel~\cref{chap:datenvorbereitung} voraus. Mathematische Details (Backpropagation) werden
konzeptuell behandelt, nicht formal bewiesen.
\end{myremark}

\section{Biologische Inspiration und das Perceptron}

\begin{mydefinition}[Perceptron]
Das \textbf{Perceptron} ist das einfachste künstliche Neuron.
Es berechnet eine gewichtete Summe seiner Eingaben und gibt 0 oder 1 aus:
\[
\hat{y} = f\!\left(\sum_{i=1}^{n} w_i x_i + b\right)
\]
wobei $f$ eine \textbf{Aktivierungsfunktion} ist (z.\,B. Schwellenwert-Funktion).
\end{mydefinition}

\begin{myexample}
Ein einzelnes Perceptron kann \enquote{UND} oder \enquote{ODER} lernen,
aber \textbf{nicht} XOR – dafür braucht es mehrere Schichten.
\end{myexample}

\section{Mehrschichtiges Netz (MLP)}

\begin{mydefinition}[Multilayer Perceptron (MLP)]
Ein \textbf{MLP} besteht aus mehreren Schichten:
\begin{itemize}
    \item \textbf{Eingabeschicht:} nimmt die Features auf
    \item \textbf{Versteckte Schichten:} lernen interne Repräsentationen
    \item \textbf{Ausgabeschicht:} gibt die Vorhersage aus
\end{itemize}
\end{mydefinition}

\begin{myexercise}[Architektur eines MLP]
Ein MLP hat 3 Eingabe-Features, zwei versteckte Schichten mit je 8 Neuronen
und eine Ausgabe.
\begin{enumerate}
    \item Zeichnen Sie (schematisch) die Netzarchitektur.
    \item Wie viele Gewichte (Parameters) hat die Verbindung von Eingabeschicht zur ersten
          versteckten Schicht?
\end{enumerate}
\begin{myanswer}
2.\ $3 \times 8 = 24$ Gewichte (plus 8 Bias-Terme = 32 Parameter).
\end{myanswer}
\end{myexercise}

\section{Aktivierungsfunktionen}

\begin{table}[H]
\centering
\small
\begin{tabular}{p{2.8cm}p{4.2cm}p{4.2cm}}
\toprule
\textbf{Name} & \textbf{Formel} & \textbf{Verwendung} \\
\midrule
Sigmoid & $\sigma(z) = \frac{1}{1+e^{-z}}$ & Ausgabe (binäre Klass.) \\
ReLU & $\text{ReLU}(z) = \max(0, z)$ & Versteckte Schichten \\
Softmax & $\text{Softmax}(z_k) = \frac{e^{z_k}}{\sum_j e^{z_j}}$ & Ausgabe (Mehrklassen) \\
\bottomrule
\end{tabular}
\caption{Häufige Aktivierungsfunktionen}
\label{tab:aktivierung-15}
\end{table}

\section{Lernen: Backpropagation (konzeptuell)}

\begin{mydefinition}[Backpropagation]
Das Netz lernt durch \textbf{Gradientenabstieg}:
\begin{enumerate}
    \item Vorhersage berechnen (\emph{Forward Pass})
    \item Fehler (Loss) messen
    \item Fehler rückwärts durch das Netz propagieren (\emph{Backpropagation})
    \item Gewichte in Richtung kleinerer Fehler anpassen
\end{enumerate}
Dieser Prozess wiederholt sich über viele \textbf{Epochen}.
\end{mydefinition}

\begin{myexercise}[Backpropagation intuitiv]
Erklären Sie mit eigenen Worten:
\begin{enumerate}
    \item Warum müssen die Gewichte zuerst zufällig initialisiert werden?
    \item Was passiert, wenn die \textbf{Lernrate} (learning rate) zu gross gewählt wird?
    \item Was bedeutet es, wenn der Loss nach jeder Epoche sinkt?
\end{enumerate}
\end{myexercise}

\section{Deep Learning}

\begin{myremark}
\textbf{Deep Learning} bezeichnet neuronale Netze mit \emph{vielen} versteckten Schichten.
Sie sind besonders leistungsfähig bei Bildern, Sprache und Texten,
benötigen aber \textbf{viele Daten} und \textbf{hohe Rechenleistung}.
\end{myremark}

\section{Grenzen: Black Box, Datenbedarf, Energie}

\begin{myexercise}[Kritische Reflexion]
Diskutieren Sie:
\begin{enumerate}
    \item Warum nennt man neuronale Netze oft eine \enquote{Black Box}?
    \item In welchen Anwendungsfeldern ist diese Intransparenz problematisch?
    \item Welche Alternative(n) würden Sie in diesen Fällen bevorzugen und warum?
\end{enumerate}
\end{myexercise}

\begin{mychallenge}[Neuronales Netz trainieren]
Trainieren Sie im Notebook ein einfaches MLP mit \texttt{sklearn.neural\_network.MLPClassifier}
auf dem Lernstil-Datensatz:
\begin{itemize}
    \item Versuchen Sie verschiedene Architekturen (\texttt{hidden\_layer\_sizes}).
    \item Vergleichen Sie die Testgenauigkeit mit dem Random Forest aus Kapitel~\cref{chap:randomforest}.
    \item Visualisieren Sie den Loss-Verlauf über die Trainingsepochen.
\end{itemize}
\end{mychallenge}

\begin{myremark}[Notebook]
Praktische Umsetzung:
\href{\notebookbaseurl06_Neuronale_Netze.ipynb}{\texttt{06\_Neuronale\_Netze.ipynb}}
\end{myremark}
